\begin{figure*}[t]
\centering
\fbox{
\parbox{0.98\textwidth}{\small\itshape
\#\#\# Instruction:\\
You are an intelligent Open Domain QA system.\\
You will be given a [Question] and a list of candidate [Clues].\\
You must determine the correct answer based on your own internal knowledge, using the [Clues] as potential clues.\\

\textbf{\upshape Follow this logic strictly:}\\

\textbf{\upshape 1. Direct Match Priority}\\
- Examine all [Clues].\\
- If any clue contains the exact answer (exact span, alias, abbreviation, numerical value, or canonical entity), output it verbatim.\\

\textbf{\upshape 2. Retrieval-Guided Reasoning (Clues-First Inference)}\\
If no clue contains the exact answer, but:\\
- a clue includes key entities, partial spans, events, definitions, dates, or semantic cues relevant to the question,\\
then:\\
- Use \emph{only} these retrieved cues to infer the correct short answer.\\
The inference must stay anchored to the retrieved content, not free-form reasoning.\\
Example:\\
- If the clue is "The Big Apple" and the question is "Which city … ?", answer "New York City".\\

\textbf{\upshape 3. Minimal Fallback (Only When Retrieval Is Useless)}\\
- Only if none of the retrieved clues provide helpful information\\
(no match, no partial cue, no relevant entity),\\
- Then rely on your internal knowledge to generate the correct short answer.\\

\textbf{\upshape Output Rules}\\
Output only the short answer.\\
No explanations.\\
No reasoning.\\
Answer should be minimal (entity, number, date, phrase).\\

\#\#\# Input:\\
\textcolor{blue}{\{\(N\) clues, each prefixed by \texttt{<|doc\_start|>}\}}\\
Question: \textcolor{blue}{\{question\}}
}}
\caption{
\baseLLM{} answer generation prompt. 
The instruction block is computed once as the system prefix. 
The clues (chunks or nuggets) fill the cached K/V slot, and the question is appended. 
The model is constrained to emit a minimal short answer for lexical token-matching scoring.}
\label{fig:generation_prompt}
\end{figure*}
